\documentclass[11pt]{article}

\usepackage[a4paper,margin=27mm]{geometry}
\usepackage{amsmath,amssymb}
\usepackage{booktabs}
\usepackage[hidelinks]{hyperref}

\newcommand{\Edr}{\mathcal{E}_{\mathrm{DR}}}
\newcommand{\norm}[1]{\left\lVert #1\right\rVert}

\title{Douglas--Rachford Line Search and Lipschitz Scaling\\
for Block-Metric Distributed Bundle Adjustment}
\author{}
\date{27 July 2026}

\begin{document}
\maketitle

\section{Purpose}

This note records what the Douglas--Rachford line-search method of Themelis,
Stella, and Patrinos actually does, how it differs from the current bundle
adjustment (BA) safeguard, and why an independent Lipschitz multiplier is a
more principled next experiment than increasing only the weak-direction
coupling term.

The relevant sources are the paper
\emph{Douglas--Rachford Splitting and ADMM for Nonconvex Optimization:
Accelerated and Newton-Type Linesearch Algorithms} and its reference
implementation in
\href{https://github.com/JuliaFirstOrder/ProximalAlgorithms.jl}
{JuliaFirstOrder/ProximalAlgorithms.jl}, specifically the \texttt{DRLS}
algorithm.

\section{Fixed-step Douglas--Rachford splitting}

For a composite objective
\[
  \min_s\;\varphi_1(s)+\varphi_2(s),
\]
one relaxed DRS step is
\begin{align}
  u_k &\in \operatorname{prox}_{\gamma\varphi_1}(s_k), \\
  v_k &\in \operatorname{prox}_{\gamma\varphi_2}(2u_k-s_k), \\
  \bar s_k &= s_k+\lambda(v_k-u_k),
  \qquad \lambda\in(0,2).
\end{align}
The fixed-point residual is
\[
  r_k=u_k-v_k,
  \qquad
  \bar s_k=s_k-\lambda r_k.
\]

For a smooth, possibly nonconvex first term with Lipschitz gradient constant
$L_f$, the reference implementation chooses
\begin{equation}
  \gamma
  =\alpha\frac{2-\lambda}{2L_f},
  \qquad \alpha=0.95 \text{ by default}.
  \label{eq:gamma}
\end{equation}
In particular, for $\lambda=1$,
\begin{equation}
  \gamma L_f=0.475.
  \label{eq:default-margin}
\end{equation}
This is a fixed stepsize selected from the smoothness bound. The official DRS
line search does not repeatedly modify $\gamma$ or the proximal coupling after
a failed accelerated direction.

\section{The official DRE line search}

Let $d_k$ be an arbitrary direction, such as an L-BFGS, Broyden, Anderson, or
Nesterov direction. The line-search center is
\begin{equation}
  s_k(\tau)
  =\tau(s_k+d_k)+(1-\tau)\bar s_k,
  \qquad \tau\in[0,1].
  \label{eq:tau-interpolation}
\end{equation}
Thus $\tau=1$ evaluates the proposed accelerated direction and $\tau=0$
evaluates the nominal DRS fallback $\bar s_k$. Backtracking halves $\tau$ and
re-evaluates the complete DRS oracle.

For the nonconvex case, define
\begin{align}
  a &= \gamma L_f, \\
  C &= \frac{\lambda}{(1+a)^2}
       \left(\frac{2-\lambda}{2}-a\right), \\
  c &= \beta C,
  \qquad \beta=0.5 \text{ by default}.
  \label{eq:decrease-constant}
\end{align}
The trial is accepted when
\begin{equation}
  \Edr\bigl(s_k(\tau)\bigr)
  \le
  \Edr(s_k)-\frac{c}{\gamma}\norm{r_k}^2.
  \label{eq:drls-acceptance}
\end{equation}
Under the exact fixed-metric assumptions, the nominal $\tau=0$ DRS fallback
satisfies the sufficient-decrease inequality. Global convergence is therefore
obtained from the fallback, while the direction is responsible only for
acceleration.

With the defaults $\lambda=1$, $\alpha=0.95$, and $\beta=0.5$,
\begin{align}
  a &= 0.475, \\
  C &= \frac{0.5-0.475}{(1+0.475)^2}
     \approx 0.01149095, \\
  c &\approx 0.00574548.
\end{align}

\section{Block-metric interpretation}

The BA implementation uses a block metric $M\succ0$ and writes the local
proximal problem as
\begin{equation}
  \min_u\;F(u)+\frac{1}{2}\norm{u-s}_M^2.
  \label{eq:block-prox}
\end{equation}
Here the scalar DRS stepsize is absorbed into $M$, so the implementation uses
$\gamma=1$ in transformed coordinates. The relevant smoothness constant is
not the inner Schur-solver constant $0.9$, but
\begin{equation}
  L_M
  =\lambda_{\max}\left(
    M^{-1/2}\nabla^2F\,M^{-1/2}
  \right).
  \label{eq:transformed-lipschitz}
\end{equation}
Equation~\eqref{eq:default-margin} therefore suggests the target
\begin{equation}
  L_M\le0.475
  \qquad (\lambda=1,\ \alpha=0.95).
  \label{eq:metric-target}
\end{equation}

A shared multiplier $\ell$ on the complete metric,
\[
  M=\ell\widehat M,
\]
is the metric analog of reducing $\gamma$. It changes proximal strength and
the DRE norm. If the same multiplier is used by every active copy, it cancels
from the block-metric consensus projection itself.

\section{Why increasing only \texttt{be} is insufficient}

The current worker uses the approximate camera metric
\begin{equation}
  M(\mathrm{be})
  =a(\mathrm{be})J^\top J
   +10\,\mathrm{be}\,\operatorname{blockdiag}(J^\top J),
  \label{eq:current-metric}
\end{equation}
where
\begin{equation}
  a(\mathrm{be})
  =\min\left(
    1.005,
    0.1\sqrt{\frac{\mathrm{be}}{\mathrm{be}_0}}
  \right).
\end{equation}
The parameter \texttt{be} consequently performs two jobs: it changes the
curvature coefficient and increases the weak-direction diagonal floor.

Under the optimistic Gauss--Newton Schur bound
\[
  S\preceq J^\top J,
\]
ignoring the additional diagonal term gives the crude estimate
\begin{equation}
  L_M\lesssim\frac{1}{a(\mathrm{be})}.
  \label{eq:crude-bound}
\end{equation}
Even at the \texttt{be=0.5} ceiling, $a=1.005$, hence
\[
  L_M\lesssim0.995,
\]
which does not reach the target $0.475$. To reach that target using only a
curvature multiple would require approximately
\begin{equation}
  M\succeq\frac{1}{0.475}S
  \approx2.105S.
  \label{eq:target-curvature-multiple}
\end{equation}

The historical Python prototype already separated the two roles:
\begin{equation}
  M=\mathrm{LipJ}\,J^\top J
    +\mathrm{be}\,D_{\mathrm{weak}}.
  \label{eq:historical-metric}
\end{equation}
Its comments report that \texttt{LipJ=2} appeared relatively safe and that
the Lipschitz multiplier should be increased after DRE failure. The prototype's
directional descent-lemma estimator was later disabled because it stalled or
failed on some scenes, so that old estimator should not be restored unchanged.

\section{Scene-931 diagnostic}

The current relative safeguard rejected iteration 7 with the following values:
\begin{center}
\begin{tabular}{lrr}
\toprule
Quantity & Iteration 6 & Iteration 7 trial \\
\midrule
Model DRE & $556{,}783$ & $548{,}423$ \\
Pixel $f(v)$ & $661{,}402$ & $984{,}722$ \\
Metric residual squared & $20{,}296.5$ & $7{,}237.95$ \\
\bottomrule
\end{tabular}
\end{center}

Using the official default $c\approx0.00574548$ and $\gamma=1$, the required
decrease based on the iteration-6 residual is only
\[
  c\norm{r_6}_M^2
  \approx0.00574548\times20{,}296.5
  \approx116.6.
\]
The observed model-DRE decrease is approximately
\[
  556{,}783-548{,}423=8{,}360,
\]
so this trial passes the formal DRLS decrease test by a large margin. It is
rejected by the current implementation because it uses
\begin{equation}
  \max\bigl(\Edr^{\mathrm{model}},f(v)\bigr)
  \label{eq:sandwich}
\end{equation}
and also checks the physical pixel objective.

The physical check cannot yet be removed. In a 90-iteration probe with only
nonfinite rejection, model-DRE decrease coexisted with pixel-SSE excursions
between $10^{16}$ and $10^{21}$, and the best SSE was only about $6.41$ million.
This does not contradict the DRLS theorem: the current BA oracle uses one
finite Gauss--Newton step, refreshes its metric, and has a special
preconditioning-only first call. It is not the exact fixed-metric proximal
oracle assumed by the theorem.

\section{Recommended controlled experiment}

Separate the curvature multiplier from the weak-direction floor:
\begin{equation}
  M=\ell J^\top J+\mathrm{be}\,D_{\mathrm{weak}}.
  \label{eq:proposed-metric}
\end{equation}
Keep \texttt{be=5e-5} fixed and sweep
\begin{equation}
  \ell\in\{1.005,\ 2.0,\ 2.105,\ 4.0\}
  \label{eq:lip-sweep}
\end{equation}
under both the DABA and DRS local trust-region policies. The value $2.105$ is
the theory-motivated threshold from
Equation~\eqref{eq:target-curvature-multiple}; $2.0$ reproduces the strongest
historical hint; and $4.0$ tests margin beyond the approximate Hessian bound.

The first experiment should use a fixed $\ell$ for the entire run. If this
removes the stale rejection sequence without unacceptable objective or runtime
regression, a second stage can estimate
Equation~\eqref{eq:transformed-lipschitz} from the reduced Schur operator and
select $\ell$ automatically. Acceleration and its $\tau$ line search should be
added only after the nominal fixed-metric oracle is stable.

\section{Conclusions}

The official DRLS method supports three conclusions relevant to the current BA
failure:
\begin{enumerate}
  \item The Lipschitz constant determines the DRS stepsize or, equivalently,
        the global scale of the proximal metric.
  \item The accelerated line search backtracks the direction toward a nominal
        DRS fallback; it does not repair a bad nominal map by increasing a
        weak-direction coupling term.
  \item The current model-DRE and physical-primal disagreement is evidence that
        finite-prox and variable-metric errors must be measured, not a reason
        to discard the physical safeguard immediately.
\end{enumerate}

The clean next change is therefore an independent fixed curvature/Lipschitz
multiplier, with \texttt{be} retained only as a numerical floor.

\end{document}